% Nouvelle version de l'introduction – réécrite avec références systématiques
\Chapter{INTRODUCTION}\label{sec:Introduction}

\section{Contexte et mise en situation}\label{sec:Contexte}

Les catastrophes naturelles — séismes, inondations, glissements de terrain, feux de forêt, ouragans — comptent parmi les menaces les plus graves pour les populations et les infrastructures à l'échelle mondiale, tant par leur fréquence croissante que par l'ampleur des dommages humains et économiques qu'elles provoquent\cite{undrr2020,noauthor_disaster_nodate}. Le rapport \emph{Human Cost of Disasters} du Bureau des Nations Unies pour la réduction des risques de catastrophes (UNDRR) recense 7\,348 événements catastrophiques majeurs entre 2000 et 2019, ayant affecté 4{,}2 milliards de personnes, causé 1{,}23 million de décès et engendré des pertes économiques directes estimées à 2\,970 milliards de dollars américains\cite{undrr2020}. Des analyses fondées sur les données EM-DAT indiquent que le nombre annuel de catastrophes reste élevé et que la variabilité interannuelle masque une tendance de fond à l'augmentation des pertes économiques, avec plusieurs années dépassant 200 milliards de dollars de dommages directs\cite{noauthor_number_nodate,burgueno_salas_global_nodate}.

Les données consolidées d'EM-DAT pour 2023--2024 confirment que les événements extrêmes continuent d'affecter des dizaines à des centaines de millions de personnes par an, avec des pertes économiques agrégées de plusieurs centaines de milliards de dollars, notamment en raison des tempêtes, inondations et vagues de chaleur dans les régions densément peuplées\cite{noauthor_number_nodate,burgueno_salas_global_nodate}. Les grands réassureurs tels que Swiss Re et Munich Re rapportent eux aussi une tendance à des pertes annuelles de plus en plus fréquentes dans la fourchette 200--300 milliards de dollars, avec une part croissante de pertes non assurées dans les pays à revenu faible et intermédiaire\cite{swiss_re_2025,munich_re_2025}. Dans ce contexte, le Cadre de Sendai pour la réduction des risques de catastrophe 2015--2030 place la compréhension des risques, le renforcement de la gouvernance et le développement de systèmes d'alerte précoce multi-aléas au cœur des priorités internationales\cite{UNDRR2015Sendai}.

Sur le plan climatique, le Sixième Rapport d'évaluation du GIEC (AR6) confirme que chaque incrément de réchauffement s'accompagne d'une augmentation de la fréquence et de l'intensité des événements extrêmes, notamment des précipitations intenses et des vagues de chaleur\cite{IPCC2021AR6}. Sous un scénario de réchauffement global de 1{,}5\,\textdegree C, l'intensité moyenne des précipitations extrêmes augmente d'environ 10\% et leur fréquence d'environ 30\% par rapport à l'ère préindustrielle\cite{IPCC2021AR6}. Les mises à jour récentes de l'Organisation météorologique mondiale confirment que la période 2015--2025 constitue la séquence de onze années les plus chaudes jamais observées, avec un réchauffement moyen annuel de l'ordre de 1{,}4--1{,}5\,\textdegree C au-dessus des niveaux préindustriels\cite{wmo2025climate}. Ces tendances renforcent l'urgence de déployer des infrastructures de surveillance résilientes et natives en intelligence artificielle pour détecter précocement les signaux annonciateurs de catastrophes et appuyer la coordination des secours.

Les réseaux de capteurs sans fil (\emph{Wireless Sensor Networks}, WSN) constituent un outil de choix pour cette surveillance fine des environnements à risque\cite{al_qundus_wireless_2020,yarali_wireless_2020}. Un WSN typique regroupe des dizaines à des centaines de nœuds capteurs miniaturisés alimentés par batterie, équipés de microcontrôleurs peu gourmands en énergie et de radios à courte portée. Ces nœuds coopèrent pour mesurer des grandeurs physiques, traiter localement les données et les acheminer par sauts multiples vers une station de base connectée à l'infrastructure d'information\cite{chong2003sensor,erdelj2017disaster}. De nombreux déploiements pilotes ont démontré la pertinence de ces architectures pour la détection précoce des feux de forêt, la surveillance de crues et la gestion des risques dans les zones montagneuses\cite{Aslan2012ForestFireWSN,biabani_energy-efficient_2020,dampage_forest_2022-1}. Ces déploiements ont aussi mis en évidence la fragilité des WSN en conditions extrêmes, avec des taux de mortalité élevés des nœuds et une dégradation rapide de la connectivité lorsque des nœuds critiques sont détruits ou défaillants\cite{AbuGhazaleh2007WSNSurvivability,erdelj2017disaster}.

Les contraintes matérielles des nœuds capteurs — batteries non rechargeables avec des budgets énergétiques de l'ordre de 0{,}5\,J, quelques dizaines à centaines de kilooctets de mémoire vive, puissances de calcul limitées — imposent une conception extrêmement parcimonieuse des protocoles de communication\cite{heinzelman2000leach,heinzelman2002leach,dietrich2009lifetime}. Dans les scénarios de surveillance de catastrophes, ces contraintes sont aggravées par la possibilité de destructions massives de nœuds, de perturbations de la topologie et d'accès humain restreint, ce qui interdit pratiquement toute opération de maintenance ou de remplacement de batterie à grande échelle\cite{erdelj2017disaster,alsayyed2025disaster}. Ces observations motivent la recherche de protocoles de routage capables de prolonger significativement la durée de vie du réseau, de maximiser la résilience topologique face aux destructions et de maintenir une qualité de service suffisante pour la détection et l'alerte.

Parallèlement, l'avènement des réseaux mobiles de sixième génération (6G) promet de transformer ces WSN en plateformes de détection et de communication natives en intelligence artificielle. Les visions 6G convergent vers des systèmes intégrant des fréquences sub-térahertz, des surfaces intelligentes reconfigurables (RIS) et des mécanismes de détection--communication intégrée (ISAC), ainsi que des architectures radio ouvertes de type O-RAN qui permettent l'exécution d'applications intelligentes sur des contrôleurs logiciels à budgets de latence stricts\cite{saad2020vision6g,polese2023oran,oran_official_architecture,etri_6g_disaster2024}. Les réseaux de capteurs pour la surveillance des catastrophes représentent un cas d'usage emblématique de cette convergence, mais posent aussi des défis inédits en termes d'efficacité énergétique, de durabilité carbone et de sécurité des décisions prises par des politiques apprises.

\section{Définitions des concepts clés}\label{sec:intro-definitions}

Cette section introduit les principaux concepts mobilisés dans la thèse afin d'en faciliter la lecture pour des publics provenant de disciplines variées.

\subsection*{Réseau de capteurs sans fil (WSN)}

Un \textbf{réseau de capteurs sans fil} (\emph{Wireless Sensor Network}, WSN) est un ensemble de petits dispositifs électroniques autonomes — les \emph{nœuds capteurs} — capables de mesurer des grandeurs physiques, de traiter localement ces mesures et de les transmettre sans fil vers une station de base au moyen de communications multi-sauts\cite{chong2003sensor,yarali_wireless_2020}. Chaque nœud est généralement alimenté par une batterie non rechargeable, dispose d'une capacité de calcul limitée (microcontrôleurs de classe ARM Cortex-M, quelques dizaines à centaines de kilooctets de RAM) et d'une radio à faible consommation énergétique\cite{yarali_wireless_2020,dietrich2009lifetime}. La durée de vie du réseau dépend fortement des décisions de communication~: la transmission de données domine largement la consommation énergétique totale, de sorte que la définition de stratégies de routage économes en énergie est un enjeu central\cite{heinzelman2000leach,dietrich2009lifetime}.

\subsection*{Routage multi-sauts}

Le \textbf{routage} désigne le processus par lequel les données produites par les capteurs sont acheminées jusqu'à la station de base au travers d'un ensemble de relais intermédiaires. Dans un WSN, la transmission directe sur de longues distances est prohibitive sur le plan énergétique, car l'énergie radio augmente typiquement avec le carré, voire la puissance quatre de la distance au-delà d'un seuil critique\cite{heinzelman2000leach,heinzelman2002leach}. Le \emph{routage multi-sauts} consiste donc à fragmenter le trajet en une succession de liens courts, souvent organisés en grappes structurées autour de \emph{chefs de grappe} (\emph{cluster heads}) qui agrègent et retransmettent les données des nœuds voisins\cite{younis2004heed,chong2003sensor}. Le choix dynamique de ces chefs de grappe et des routes de transfert — « qui transmet quoi, vers qui, et quand » — constitue le cœur du problème d'optimisation abordé dans cette thèse\cite{metaheuristic_wsn_slr2026,chaurasia_mocraw_2023}.

\subsection*{Apprentissage par renforcement multi-agent (MARL)}

L'\textbf{apprentissage par renforcement} (\emph{Reinforcement Learning}, RL) permet à un agent d'apprendre une politique de décision à partir d'interactions avec un environnement, en maximisant une somme de récompenses cumulées\cite{mnih2015dqn}. L'\textbf{apprentissage par renforcement multi-agent} (\emph{Multi-Agent RL}, MARL) étend ce paradigme à un ensemble d'agents partageant un environnement commun et des objectifs souvent coopératifs\cite{noauthor_multi-agent_nodate}. Dans le cadre de cette thèse, chaque nœud capteur est modélisé comme un agent coopératif qui apprend à prendre des décisions de routage locales de manière à optimiser simultanément la livraison de paquets et la consommation d'énergie au niveau global du réseau\cite{rashid2018qmix,son2019qtran,Hayes2022MOMARL}. Les approches de décomposition de valeur comme QMIX et QTRAN fournissent un cadre puissant pour factoriser la valeur d'équipe en contributions individuelles tout en respectant certaines propriétés structurelles\cite{rashid2018qmix,son2019qtran,Wang2021QPLEX}.

\subsection*{Réseaux de neurones de graphes (GNN)}

Un \textbf{réseau de neurones de graphes} (\emph{Graph Neural Network}, GNN) est une famille de modèles neuronaux conçus pour traiter des données structurées sous forme de graphes, en exploitant explicitement les relations d'adjacence entre nœuds\cite{gilmer2017mpnn,kipf2017semi,velickovic2018gat}. Dans un WSN, la topologie physique et logique est naturellement représentée par un graphe\cite{al_qundus_wireless_2020,erdelj2017disaster}. En appliquant des opérations de \emph{message passing}, un GNN produit pour chaque nœud une représentation vectorielle qui agrège l'information locale (énergie résiduelle, file d'attente, position) et le contexte global (centralité, appartenance à des structures critiques), fournissant ainsi un codage riche pour guider les décisions de routage ou la sélection de politiques expertes\cite{liu2024gnncomm,Nguyen2025GNNIoT}.

\subsection*{Analyse du cycle de vie (ACV) et intensité carbone}

L'\textbf{analyse du cycle de vie} (ACV, \emph{Life Cycle Assessment}, LCA) est une méthodologie normalisée par la norme ISO~14040 pour quantifier les impacts environnementaux d'un produit ou d'un système « du berceau à la tombe »\cite{iso14040}. Des travaux récents appliqués à l'Internet des objets montrent que le \emph{carbone incorporé} lors de la fabrication des nœuds domine très largement les émissions opérationnelles sur toute la durée de vie, avec des ordres de grandeur pouvant atteindre un facteur million entre ces deux contributions\cite{pirson2021embodied,bashir2024sunk,mersy2025lca_data}. Cette thèse introduit la notion d'\textbf{intensité carbone du cycle de vie} (LCI, \emph{Lifecycle Carbon Intensity}) pour les WSN, qui exprime l'empreinte carbone totale en grammes de CO$_2$ équivalent par paquet utile livré. Cette métrique rend compte à la fois de l'efficacité de livraison des données et du coût environnemental complet de la fabrication, de l'exploitation et de la fin de vie des nœuds\cite{pirson2021embodied,aioti2024carbon}.

\subsection*{Confiance zéro (zero-trust)}

Le paradigme \textbf{zéro-confiance} (\emph{zero-trust security}) repose sur le principe « ne jamais faire confiance, toujours vérifier »\cite{rose2020nist}. Contrairement aux modèles de sécurité périmétrique classiques qui considèrent les entités internes comme implicitement fiables, les architectures zéro-confiance exigent que chaque requête, chaque flux et chaque décision soient authentifiés, autorisés et évalués en continu\cite{Kizza2013,ztaiot_acm2025}. Dans le contexte des WSN pilotés par IA, cela signifie que les décisions de routage proposées par des politiques apprises — potentiellement vulnérables à des attaques adversariales ou à des hallucinations hors distribution — doivent être interceptées, évaluées et éventuellement corrigées par un mécanisme de contrôle externe, doté de garanties formelles de sûreté\cite{amodei2016concrete,Huang2017AdvRL,subplay_marl2025}.

\subsection*{Sixième génération de réseaux mobiles (6G) et O-RAN}

La \textbf{sixième génération de réseaux mobiles} (6G) est envisagée comme une plateforme de communication et de détection native en IA, intégrant des fréquences sub-térahertz, des réseaux non terrestres (NTN), des surfaces intelligentes reconfigurables et des fonctions d'orchestration avancées pour les services critiques\cite{saad2020vision6g,amatetti2024ntn,abualigah_swarm_2023}. L'architecture \textbf{O-RAN} (\emph{Open Radio Access Network}) découple les fonctions logicielles et matérielles de la couche d'accès radio, expose des interfaces ouvertes et introduit des contrôleurs intelligents (near-RT RIC) où des xApps et rApps peuvent être déployées pour optimiser le réseau en temps quasi réel\cite{polese2023oran,oran_official_architecture,oran_survey_acm2024}. Les contraintes de latence ($\sim$ 10~ms pour le near-RT RIC) et de ressources des équipements d'extrémité imposent des exigences strictes sur la compacité et l'efficacité des algorithmes embarqués\cite{dasilva2024xapp,demichelis2025edge,demirel2026agentic}.

\section{Justification économique, sociétale et réglementaire}\label{sec:intro-justification}

Les enjeux économiques et sociétaux associés aux catastrophes naturelles justifient la recherche de solutions de surveillance plus efficaces et plus résilientes. Les rapports annuels d'EM-DAT et des grands réassureurs indiquent que les pertes économiques directes liées aux catastrophes se situent régulièrement dans la fourchette des centaines de milliards de dollars par an\cite{undrr2020,burgueno_salas_global_nodate,swiss_re_2025}. Le Cadre de Sendai estime qu'un dollar investi dans la prévention et les systèmes d'alerte précoce peut générer un retour de 4 à 10 dollars en pertes évitées, illustrant la forte valeur socioéconomique des infrastructures de surveillance performantes\cite{UNDRR2015Sendai}. L'indicateur C-2 du cadre de Sendai cible spécifiquement le nombre de personnes couvertes par des systèmes d'alerte précoce multi-aléas, renforçant la pertinence des WSN intelligentés étudiés dans cette thèse\cite{UNDRR2015Sendai,wmo2025climate}. Ces bénéfices sont particulièrement critiques dans les pays à revenu faible et intermédiaire, qui concentrent la majorité des décès liés aux catastrophes\cite{undrr2020,wmo2025climate}.

En parallèle, la transition numérique et le développement accéléré de l'intelligence artificielle ont placé l'empreinte environnementale des TIC au centre des débats. Des études récentes estiment la contribution actuelle du secteur des TIC aux émissions mondiales de gaz à effet de serre entre 2 et 4~\%\cite{freitag2021ict,andrae2021real_climate}. Dans ce contexte, les WSN pour la surveillance des catastrophes doivent être conçus de manière à maximiser la valeur environnementale nette~: prolonger la durée de vie des nœuds, amortir le carbone incorporé sur un nombre maximal de paquets utiles et exploiter les signaux d'intensité carbone du réseau électrique lorsque cela est pertinent\cite{pirson2021embodied,aioti2024carbon,carbon_signals_acm2025}.

Sur le plan réglementaire, plusieurs cadres structurent désormais explicitement les exigences en matière de durabilité et de sécurité des systèmes d'IA. Outre le Cadre de Sendai\cite{UNDRR2015Sendai} et la norme ISO~14040\cite{iso14040}, le \emph{NIST AI Risk Management Framework} propose une approche systématique pour identifier, évaluer et atténuer les risques liés aux systèmes d'IA\cite{nist2023airisk}. Au niveau européen, l'\emph{Artificial Intelligence Act} introduit une catégorisation des systèmes d'IA par niveau de risque, dont font partie les applications critiques en gestion de catastrophes\cite{euaiact2024}. Le cadre zéro-confiance du NIST SP~800-207 et les travaux récents sur la sécurité des architectures O-RAN soulignent la nécessité de mécanismes de contrôle de décision explicites pour les politiques apprises\cite{rose2020nist,polese2023oran,ahmadi2025zt_segmentation,ztaiot_acm2025}.

La pile proposée dans cette thèse — combinant un routage écoénergétique, une minimisation de l'intensité carbone du cycle de vie et un contrôle de décision zéro-confiance — vise à répondre simultanément à des impératifs économiques, environnementaux et réglementaires\cite{freitag2021ict,pirson2021embodied,rose2020nist,itu_ict2025}.

\section{Éléments de la problématique}\label{sec:Problematique}

L'analyse de la littérature détaillée au Chapitre~\ref{sec:RevLitt} met en évidence trois défis fondamentaux qui restent insuffisamment adressés dans le contexte des WSN de surveillance des catastrophes intégrés aux technologies 6G~: (1)~l'efficacité énergétique du routage, (2)~la durabilité carbone du cycle de vie, et (3)~la sécurité des politiques de routage apprises par IA\cite{erdelj2017disaster,freitag2021ict,metaheuristic_wsn_slr2026,alsayyed2025disaster}.

\subsection{Défi 1 : Efficacité énergétique du routage dans les WSN de surveillance des catastrophes}

Le premier défi concerne le maintien d'une \emph{durée de vie réseau} suffisante dans des environnements hostiles, où le remplacement des batteries ou des nœuds est pratiquement impossible. Les protocoles de routage hiérarchiques classiques comme LEACH, HEED et leurs variantes reposent sur des heuristiques statiques pour l'élection des chefs de grappe et n'exploitent que partiellement l'information topologique globale\cite{heinzelman2000leach,heinzelman2002leach,younis2004heed,kumar2012eeleach}. Des approches plus récentes basées sur des métaheuristiques améliorent la consommation énergétique et la durée de vie, mais au prix de temps de convergence élevés et sans garanties fortes en environnements non stationnaires\cite{karaboga2005abc,abualigah_swarm_2023,chaurasia_mocraw_2023,elkhediri2024abcaco,dinesh_hbo-sroa_2024}. Une revue systématique couvrant la période 2019--2024 constate que ces approches se concentrent essentiellement sur l'énergie et la durée de vie, sans intégrer explicitement la dimension carbone ni les contraintes d'embarqué TinyML\cite{metaheuristic_wsn_slr2026,tinyml_review_acm2024}.

Le MARL a récemment été identifié comme une alternative prometteuse pour le routage dans les réseaux dynamiques, en particulier grâce aux méthodes de décomposition de valeur telles que QMIX, QTRAN et leurs extensions\cite{rashid2018qmix,son2019qtran,liu2024gnncomm,ding_packet_2022}. Toutefois, ces approches souffrent de plusieurs limitations critiques lorsqu'on envisage leur déploiement sur des nœuds capteurs contraints~: (i)~une explosion mémoire liée à l'utilisation de réseaux récurrents et de fonctions de valeur d'équipe de grande dimension, (ii)~une prise en compte limitée de la structure de graphe sous-jacente au réseau, et (iii)~l'absence de cadres de \emph{fusion dynamique} de politiques expertes complémentaires en fonction du contexte topologique local\cite{liu2024gnncomm,guo2024rldas,Hayes2022MOMARL}. Il manque ainsi un cadre de routage qui combine la puissance de plusieurs politiques MARL expertes avec un encodeur GNN compact, de manière compatible avec les contraintes de mémoire et de calcul des plates-formes capteurs\cite{Nguyen2025GNNIoT,tinyml_review_acm2024}.

\subsection{Défi 2 : Empreinte carbone du cycle de vie des WSN}

Le deuxième défi concerne la \emph{durabilité carbone} du cycle de vie des WSN. Pour les dispositifs IoT à faible consommation, des analyses d'ACV détaillées révèlent que le carbone incorporé lors de la fabrication des nœuds domine de plusieurs ordres de grandeur le carbone opérationnel généré par leur consommation électrique\cite{pirson2021embodied,bashir2024sunk,aioti2024carbon}.

Ce constat a deux implications majeures~: d'une part, réduire la consommation énergétique opérationnelle d'un capteur ne diminue que marginalement son empreinte carbone totale si le nombre de paquets utiles livrés reste faible~; d'autre part, toute destruction prématurée de nœuds se traduit par un gaspillage complet du carbone incorporé\cite{pirson2021embodied,mersy2025lca_data}. Par ailleurs, l'intensité carbone du réseau électrique varie fortement selon la localisation et le moment, ce qui ouvre la voie à des stratégies de modulation de la charge computationnelle en fonction des signaux d'intensité carbone disponibles\cite{ember2024,uk_carbon_api,carbon_signals_acm2025}. Aucun protocole de routage existant ne prend explicitement en compte l'empreinte carbone du cycle de vie dans ses décisions, notamment dans le contexte des WSN de surveillance des catastrophes\cite{freitag2021ict,aioti2024carbon,metaheuristic_wsn_slr2026}.

\subsection{Défi 3 : Sécurité des politiques de routage apprises par intelligence artificielle}

Le troisième défi concerne la \emph{sûreté et la sécurité} des politiques de routage apprises par IA. L'introduction de l'IA dans la boucle de contrôle des réseaux élargit la surface d'attaque et introduit de nouvelles formes de défaillances, notamment en situation hors distribution\cite{amodei2016concrete,garcia2015saferl,brunke2022safe}. Des travaux récents ont montré que des politiques de RL peuvent être vulnérables à des attaques adversariales, des backdoors ou des perturbations ciblées\cite{Huang2017AdvRL,Gleave2020AdvPolicies,Kiourti2021BackdoorDRL,subplay_marl2025,ami_attack_marl2024}. En particulier, des politiques dites \emph{voyous} peuvent dégrader délibérément les performances de systèmes multi-agents coopératifs --- justifiant la nécessité d'un mécanisme de contrôle externe capable d'intercepter et de neutraliser ces comportements déviants.

Les cadres de \emph{safe RL} existants, tels que l'optimisation de politiques sous contraintes (CPO) ou le \emph{shielding} réactif, fournissent des garanties de sûreté en limitant l'exploration à des régions jugées sûres ou en filtrant les actions non valides\cite{achiam2017cpo,alshiekh2018shielding,garcia2015saferl}. Toutefois, ces approches présentent généralement deux limitations majeures pour les WSN~: (i)~elles opèrent souvent avec une granularité binaire (autoriser/bloquer), ce qui peut dégrader la connectivité lorsqu'elles sont appliquées de manière conservatrice, et (ii)~elles ne s'intègrent pas explicitement dans un cadre zéro-confiance aligné sur les recommandations du NIST et des architectures 6G/O-RAN\cite{rose2020nist,polese2023oran,pranger2025shields}. Aucun travail existant ne réunit simultanément un encodeur GNN, un évaluateur de confiance-risque neuronal et un \emph{bouclier déterministe à décisions graduées} pour contrôler en temps réel les décisions de routage apprises tout en respectant les budgets de latence des architectures O-RAN\cite{corsi2024verification,dmps_shielding2024,ztaiot_acm2025}.

\section{Questions de recherche}\label{sec:Questions}

Face à ces trois défis, la thèse s'articule autour d'une question de recherche principale et de trois questions spécifiques correspondant à chacun des articles qui composent la pile GAPF--CALASH--CASTER-ZT.

\medskip
\begin{quote}
\textbf{Question principale :} \emph{Comment concevoir une pile complète de protocoles de routage intelligent — écoénergétique, durable en carbone et sécurisée — pour les réseaux de capteurs de surveillance des catastrophes natifs en intelligence artificielle et intégrés aux technologies 6G ?}
\end{quote}

\medskip
Cette question principale se décline en trois questions spécifiques~:

\begin{description}
  \item[Q1 — Efficacité énergétique (Article~1)]~:\emph{Est-il possible d'atteindre un routage multi-sauts quasi optimal dans les WSN de surveillance des catastrophes en fusionnant des politiques MARL expertes complémentaires au moyen d'un encodeur GNN compact, tout en respectant les contraintes de mémoire et de calcul des plates-formes embarquées ?}\cite{gapf_tgcn2026,rashid2018qmix,son2019qtran,liu2024gnncomm,tinyml_review_acm2024}

  \item[Q2 — Durabilité carbone (Article~2)]~:\emph{Peut-on réduire significativement l'intensité carbone du cycle de vie (LCI) d'un WSN de surveillance des catastrophes en intégrant conjointement la réduction de données sensible au carbone, le routage sous contraintes de budget carbone, l'auto-guérison autonomique et la comptabilité ACV, le tout sur une couche physique 6G ?}\cite{pirson2021embodied,freitag2021ict,aioti2024carbon,itu_ict2025}

  \item[Q3 — Sécurité zéro-confiance (Article~3)]~:\emph{Peut-on garantir la détection et la mitigation des politiques de routage compromises ou erronées avec un taux de faux positifs nul ou négligeable, en couplant un encodeur GNN, un évaluateur de confiance-risque et un bouclier déterministe à décisions graduées, tout en respectant le budget de latence du near-RT RIC O-RAN ?}\cite{rose2020nist,polese2023oran,pranger2025shields,ztaiot_acm2025}
\end{description}

\section{Objectifs de recherche}\label{sec:Objectifs}

\subsection{Objectif général}

L'objectif général de la thèse est de \textbf{concevoir, implémenter et évaluer une pile complète de protocoles de routage intelligent — écoénergétique, durable en carbone et sécurisée — pour les réseaux de capteurs de surveillance des catastrophes intégrés aux technologies 6G}, en apportant des contributions originales sur les plans algorithmique, méthodologique et expérimental\cite{freitag2021ict,saad2020vision6g,polese2023oran}.

\subsection{Objectifs spécifiques}

Cet objectif général se décline en trois objectifs spécifiques correspondant aux trois articles principaux~:

\begin{enumerate}
  \item \textbf{OS1 — Fusion légère de politiques MARL pour le routage écoénergétique (Article~1 : GAPF)}\cite{gapf_tgcn2026,rashid2018qmix,son2019qtran,liu2024gnncomm,Nguyen2025GNNIoT} :
  \begin{itemize}
    \item concevoir un encodeur GCN peu profond produisant des plongements de nœuds compacts exploitables pour la sélection de politiques ;
    \item développer un contrôleur de type \emph{Contextual Algorithm Selection} capable de fusionner dynamiquement plusieurs politiques MARL expertes en fonction du contexte topologique ;
    \item garantir un compromis favorable entre taux de livraison de paquets, consommation d'énergie et empreinte mémoire pour des WSN de surveillance des catastrophes.
  \end{itemize}

  \item \textbf{OS2 — Routage conscient du carbone du cycle de vie avec auto-guérison (Article~2 : CALASH)}\cite{pirson2021embodied,aioti2024carbon,mersy2025lca_data,itu_ict2025} :
  \begin{itemize}
    \item définir et formaliser la métrique LCI pour les WSN et l'intégrer dans le problème de routage ;
    \item concevoir une architecture en quatre piliers (réduction de données sensible au carbone, contrôle de file et de puissance sous contraintes de budget carbone, auto-guérison autonomique, comptabilité ACV) ;
    \item établir des garanties théoriques de respect de budgets carbone via l'optimisation de Lyapunov et valider empiriquement les gains de LCI par rapport aux approches de référence.
  \end{itemize}

  \item \textbf{OS3 — Contrôle de décision zéro-confiance pour les réseaux natifs en IA (Article~3 : CASTER-ZT)}\cite{rose2020nist,garcia2015saferl,pranger2025shields,ztaiot_acm2025} :
  \begin{itemize}
    \item élaborer un encodeur GNN couplé à un évaluateur neuronal de confiance-risque, tenant compte de l'incertitude épistémique des modèles ;
    \item concevoir un bouclier déterministe à décisions graduées compatible avec les budgets de latence de l'O-RAN near-RT RIC ;
    \item formuler et démontrer un ensemble de propriétés formelles de sûreté et les valider expérimentalement sur des scénarios de menaces réalistes.
  \end{itemize}
\end{enumerate}

\section{Portée et délimitation de la thèse}\label{sec:Portee}

La portée de la thèse est délimitée selon quatre axes complémentaires.

\begin{enumerate}
  \item \textbf{Domaine applicatif.} Les contributions sont conçues et évaluées exclusivement dans le contexte des WSN de surveillance des catastrophes (feux de forêt, inondations, séismes et glissements de terrain), caractérisés par des contraintes énergétiques sévères, des topologies potentiellement fragmentées et des exigences de fiabilité élevées\cite{erdelj2017disaster,alsayyed2025disaster}. Les scénarios IoT génériques (domotique, industrie 4.0) ou les réseaux mobiles grand public ne constituent pas un objet d'étude direct, même si certaines techniques développées pourraient y être transférées.

  \item \textbf{Échelle et environnement.} Les protocoles sont évalués sur des réseaux de 20 à 200 nœuds capteurs statiques ou quasi statiques, ce qui correspond à la plage d'échelle typiquement étudiée dans la littérature pour les scénarios de surveillance de catastrophes\cite{erdelj2017disaster,biabani_energy-efficient_2020,dampage_forest_2022-1}. La mobilité continue des nœuds et les déploiements à très grande échelle sont laissés en dehors du périmètre et identifiés comme des directions de recherche futures.

  \item \textbf{Couche d'abstraction.} Les contributions se situent principalement aux couches réseau et transport (routage multi-sauts, gestion des files, contrôle de décision). La couche physique 6G est modélisée à un niveau d'abstraction suffisant pour évaluer l'impact des choix de routage, sans prétendre redéfinir les spécifications radio\cite{saad2020vision6g,ns3_thz,amatetti2024ntn,alhilo2024ris_uav}.

  \item \textbf{Environnement de validation.} Les résultats expérimentaux reposent sur des environnements de simulation conformes à la littérature et sur des jeux de données réels de capteurs pour la validation\cite{intel_lab_data,erdelj2017disaster,dampage_forest_2022-1}. L'évaluation sur bancs d'essai matériels réels est identifiée comme une étape ultérieure, en cohérence avec les recommandations récentes en matière de reproductibilité en apprentissage machine\cite{Pineau2021Reproducibility,demichelis2025edge,demirel2026agentic}.
\end{enumerate}

En résumé, la thèse se concentre sur la conception d'une \emph{pile protocolaire complète} — GAPF pour l'efficacité énergétique, CALASH pour la durabilité carbone et CASTER-ZT pour la sécurité zéro-confiance — adaptée aux WSN de surveillance des catastrophes natifs en IA et intégrés aux architectures 6G/O-RAN.

\section{Plan de la thèse}\label{sec:Plan}

La thèse, rédigée sous forme d'articles, est organisée en huit chapitres. Le Chapitre~\ref{sec:RevLitt} présente la revue de littérature couvrant le routage WSN, le MARL et les GNN, la durabilité carbone, la sécurité des systèmes autonomes et les technologies 6G, et en dégage les lacunes qui motivent la pile proposée\cite{chong2003sensor,freitag2021ict,saad2020vision6g,polese2023oran,metaheuristic_wsn_slr2026}. Le Chapitre~\ref{sec:Demarche} explicite la démarche méthodologique, le cadre expérimental commun aux trois articles et les liens entre les contributions, en s'appuyant sur les méthodologies de \emph{design science} et les recommandations récentes en matière de reproductibilité\cite{Hevner2004DesignScience,Peffers2007DSRM,Pineau2021Reproducibility}. Les Chapitres~\ref{sec:Theme1}, \ref{sec:Theme2} et~\ref{sec:Theme3} présentent respectivement les articles GAPF (publié dans IEEE Transactions on Green Communications and Networking, vol.~10, pp.~2921--2933, 2026, DOI:~\texttt{10.1109/TGCN.2026.3690507})\cite{gapf_tgcn2026}, CALASH et CASTER-ZT\cite{liu2024gnncomm,pirson2021embodied,pranger2025shields,Nguyen2025GNNIoT}. Le Chapitre~\ref{sec:Discussion} réalise une analyse transversale des résultats, discute les synergies et les limites de la pile proposée et met en perspective les implications pour la conception de réseaux natifs en IA critiques. Enfin, le Chapitre~\ref{sec:Conclusion} synthétise les contributions, rappelle les principales limitations et identifie des pistes de recherche futures\cite{Nguyen2024DTORAN,demirel2026agentic}.
